We first show why the additional policy restriction in the proposed statement is load-bearing.  Write
\[
 E_t=\{R_t(x_t)=y_t\},\qquad E=\bigcap_{t\in[H]}E_t,
 \qquad B=\{A_t=a_t\text{ for every }t\in[H]\}.
\]
Because the initial state is fixed at \(s_1\), the deterministic structural equations in \(\text{ass:markov-structural-equations}\) imply, by induction over \(t\), that the trajectory event is exactly \(\{\tau\}=B\cap E\).  Indeed, on \(B\), the equality \(R_t(x_t)=y_t=s_{t+1}\) propagates the prescribed state from time \(t\) to time \(t+1\), and the converse follows directly from the definition of \(\tau\).

Let \(\lambda_t\) be the law of \(R_t\).  Since \(R_t\) is a deterministic function of \(U_t\), \(\text{ass:independent-time-disturbances}\) implies that \(R_1,\ldots,R_H\) are independent.  Hence their joint mass function is \(\prod_{t\in[H]}\lambda_t\).  For any response-table tuple \(\rho=(\rho_t:t\in[H])\) in its support and satisfying \(\rho_t(x_t)=y_t\) for every \(t\), the external-randomization condition in the proposed statement gives
\[
 \Pr_{\mathcal M}(B\mid S_1=s_1,R_t=\rho_t\text{ for all }t)
 =c_\tau,
 \qquad
 c_\tau:=\prod_{t\in[H]}\pi_t(a_t\mid s_t)>0.
\]
Therefore, for every subset \(D\) of the finite product response-table space,
\[
 \Pr_{\mathcal M}\{(R_t)_{t\in[H]}\in D,\tau\}
 =c_\tau\Pr_{\mathcal M}\{(R_t)_{t\in[H]}\in D\cap E\}.
\]
By \(\text{ass:transition-marginals}\),
\[
 \Pr(E_t)=K_t(y_t\mid x_t)=p_t,
\]
and \(p_t>0\) by \(\text{ass:factual-positivity}\).  Independence across time now gives
\[
 \Pr(\tau)=c_\tau\prod_{t\in[H]}p_t>0
\]
and, after division by this positive quantity,
\[
 \begin{aligned}
 &\Pr_{\mathcal M}\{R_t=\rho_t\text{ for all }t\mid\tau\}\\
 &\qquad=\prod_{t\in[H]}
 \frac{\lambda_t(\rho_t)\mathbf 1\{\rho_t(x_t)=y_t\}}{p_t}.
 \end{aligned}
\]
Thus conditioning on \(\tau\) preserves independence of the time blocks and conditions the time-\(t\) table only on \(E_t\).

Define, for \(t\in[H]\), \(x\in\mathcal X\), and \(y\in\mathcal S\),
\[
 q_{t,x}(y)=\Pr_{\mathcal M}\{R_t(x)=y\mid E_t\}.
\]
The preceding factorization also makes this equal to \(\Pr_{\mathcal M}\{R_t(x)=y\mid\tau\}\).  On \(E_t\), \(R_t(x_t)=y_t\), so \(q_{t,x_t}=\delta_{y_t}\).  For \(x\ne x_t\), the multiplication rule, event inclusion, and \(\text{ass:transition-marginals}\) give
\[
 \begin{aligned}
 p_tq_{t,x}(y)
 &=\Pr\{R_t(x_t)=y_t,R_t(x)=y\}\\
 &\le \Pr\{R_t(x)=y\}
 =K_t(y\mid x).
 \end{aligned}
\]
Each \(q_{t,x}\) is a probability mass function because it is a conditional law on a positive-probability event.  These are exactly the constraints in \(\text{def:posterior-row-fiber}\), so \(q\in\mathcal B(K,\tau)\).

It remains to prove the converse and attainment.  Fix \(q\in\mathcal B(K,\tau)\).  At each time \(t\), draw \(Z_t\) with law \(K_t(\cdot\mid x_t)\) and set \(R_t(x_t)=Z_t\).  Suppose first that \(p_t<1\).  For every \(x\ne x_t\), define
\[
 h_{t,x}(y)=\frac{K_t(y\mid x)-p_tq_{t,x}(y)}{1-p_t},
 \qquad y\in\mathcal S.
\]
The numerator is nonnegative by \(\text{def:posterior-row-fiber}\), while
\[
 \sum_{y\in\mathcal S}h_{t,x}(y)
 =\frac{1-p_t}{1-p_t}=1.
\]
Thus \(h_{t,x}\in\Delta(\mathcal S)\).  Conditional on \(Z_t=y_t\), draw the coordinates \(R_t(x)\), \(x\ne x_t\), independently with respective laws \(q_{t,x}\).  Conditional on \(Z_t\ne y_t\), draw them independently with respective laws \(h_{t,x}\).  This defines one joint response-table law \(\lambda_t\), and for every \(x\ne x_t\),
\[
 \Pr\{R_t(x)=y\}
 =p_tq_{t,x}(y)+(1-p_t)h_{t,x}(y)
 =K_t(y\mid x).
\]
The factual coordinate has marginal \(K_t(\cdot\mid x_t)\) by construction.  Moreover,
\[
 \Pr\{R_t(x)=y\mid R_t(x_t)=y_t\}=q_{t,x}(y).
\]
If \(p_t=1\), the fiber inequality says \(q_{t,x}(y)\le K_t(y\mid x)\) for every \(y\); since both sides sum to one, equality holds coordinatewise.  In this case draw all nonfactual coordinates with laws \(q_{t,x}=K_t(\cdot\mid x)\) in the sole branch \(Z_t=y_t\).  Hence the same marginal and conditional conclusions hold in both cases.

Draw the constructed tables independently across \(t\), take \(U_t=R_t\), and define \(F_t(x,U_t)=U_t(x)\).  The lookup equation verifies \(\text{ass:markov-structural-equations}\), the independent draws verify \(\text{ass:independent-time-disturbances}\), and the preceding marginal calculation verifies \(\text{ass:transition-marginals}\).  Consequently the constructed environment lies in \(\mathfrak M(K)\) by \(\text{def:compatible-scm-class}\).  Start it at \(s_1\) and couple it to external factual policy randomization with kernel \(\pi\).  The support condition and factual positivity give \(\Pr(\tau)=c_\tau\prod_t p_t>0\); the already proved factorization then shows that its posterior row kernel is exactly the prescribed \(q\).

Finally, this construction chooses one posterior response-table law at each time and samples it once.  Every coalition intervention changes only the policy kernel used to select a row; it does not redraw or reselect the response table.  Conditional on \(\tau\), the time-\(t\) table is independent of earlier posterior time blocks, and a coalition path queries the single coordinate \(R_t(S_t,A_t)\), whose conditional marginal is \(q_{t,S_t,A_t}\).  Thus the same \(q\), and the same underlying table draw, evaluate all nonfactual rows and every coalition world.  This proves both assertions of the proposed statement.